% ===========================================================================
%  Chapitre 8 — La fonction exponentielle
%  PE 2026, composante ANALYSE
% ===========================================================================
\bmtchapitre{La fonction exponentielle}
  {Connaître la fonction exponentielle et ses propriétés algébriques ;
   résoudre les équations du type $e^{u(x)} = a$, $e^{u(x)} = e^{v(x)}$ et
   $a\,e^{2x}+b\,e^{x}+c = 0$ ; étudier une fonction contenant une
   exponentielle.}
  {Analyse \textbullet\ environ 4 heures}

L'exponentielle est le logarithme lu à l'envers. Tout ce que $\ln$ fait,
$\exp$ le défait : $\ln$ transforme les produits en sommes, $\exp$ transforme
les sommes en produits. C'est pourquoi ce chapitre est court — presque chaque
propriété est le reflet d'une propriété déjà connue.

Ce qui change, c'est le comportement : là où le logarithme monte lentement,
l'exponentielle explose. Elle décrit tout ce qui croît en se multipliant par
lui-même : une population, un capital, une épidémie. Et à l'examen, elle
partage avec le logarithme la vedette du problème.

% ---------------------------------------------------------------------------
\begin{revision}
Le chapitre 7, et rien d'autre.

\begin{enumerate}[label=\textbf{\arabic*.}, leftmargin=*, itemsep=0.35em]
  \item Que vaut $\ln 1$ ? $\ln e$ ?
  \item Résoudre $\ln x = 2$.
  \item Simplifier $\ln(a b)$ et $\ln\dfrac ab$ pour $a>0$ et $b>0$.
  \item Sur quel ensemble $\ln$ est-elle définie ? Est-elle croissante ou
        décroissante ?
  \item Résoudre $X^{2}-4X+3 = 0$.
  \item Dériver $x \mapsto \ln(2x+1)$ sur $\intervalleo{-\frac12}{+\infty}$.
\end{enumerate}
\end{revision}

\begin{automatismes}
Après le cours, sans calculatrice.

\begin{enumerate}[label=\textbf{\alph*)}, leftmargin=*, itemsep=0.25em]
  \item $e^{0}$
  \item $e^{1}$
  \item $e^{2}\times e^{3}$
  \item $\dfrac{e^{5}}{e^{2}}$
  \item $\left(e^{2}\right)^{3}$
  \item $e^{-1}$ sous forme de fraction
  \item $\ln\left(e^{7}\right)$
  \item Le signe de $e^{x}$ pour $x = -100$
\end{enumerate}
\end{automatismes}

\begin{corrige}[automatismes]
a) $1$ \quad b) $e$ \quad c) $e^{5}$ \quad d) $e^{3}$ \quad e) $e^{6}$
\quad f) $\frac1e$ \quad g) $7$ \quad h) strictement positif.
\end{corrige}

% ===========================================================================
\section{Définition et propriétés}

\begin{definition}[fonction exponentielle]
La fonction $\ln$ est strictement croissante de $\intervalleo{0}{+\infty}$
sur $\R$ : tout réel $y$ possède donc un unique antécédent $x>0$ tel que
$\ln x = y$. La fonction qui à $y$ associe cet antécédent s'appelle
\textbf{fonction exponentielle}, notée $\exp$ ou $x \mapsto e^{x}$.
\end{definition}

\begin{propriete}[les deux égalités de base]
Pour tout réel $x$ et tout réel $y>0$ :
\[
  \ln\left(e^{x}\right) = x, \qquad e^{\ln y} = y .
\]
En particulier $e^{0} = 1$ et $e^{1} = e$.
\end{propriete}

\begin{theoreme}[propriétés algébriques]
Pour tous réels $a$ et $b$, et tout entier $n$ :
\[
  e^{a+b} = e^{a}\times e^{b}, \qquad
  e^{-a} = \frac{1}{e^{a}}, \qquad
  e^{a-b} = \frac{e^{a}}{e^{b}}, \qquad
  \left(e^{a}\right)^{n} = e^{na}.
\]
\end{theoreme}

\begin{demonstration}[de la première égalité]
Posons $A = e^{a}$ et $B = e^{b}$ ; ce sont deux réels strictement positifs,
et $\ln A = a$, $\ln B = b$. Alors
$\ln(AB) = \ln A+\ln B = a+b$, donc $AB = e^{a+b}$ par définition de
l'exponentielle.
\end{demonstration}

\begin{propriete}[signe, variations, limites]
Pour tout réel $x$, $e^{x} > 0$ : \textbf{l'exponentielle ne s'annule
jamais et n'est jamais négative}. Elle est dérivable sur $\R$, de dérivée
elle-même :
\[
  \left(e^{x}\right)' = e^{x},
\]
elle est donc strictement croissante sur $\R$. Enfin
\[
  \limminf e^{x} = 0, \qquad \limpinf e^{x} = +\infty, \qquad
  \limpinf \frac{e^{x}}{x} = +\infty, \qquad \limminf x\,e^{x} = 0 .
\]
\end{propriete}

\begin{avertissement}
Les deux dernières limites sont les formes indéterminées de ce chapitre.
Elles disent la même chose sous deux angles : \emph{l'exponentielle
l'emporte toujours sur $x$}. En $+\infty$ elle est infiniment plus grande ;
en $-\infty$ elle est infiniment plus petite, au point d'écraser le facteur
$x$. Les sujets les donnent en indication, comme pour le logarithme.
\end{avertissement}

\begin{visualisation}[deux courbes symétriques]
\begin{center}
\begin{tikzpicture}
\begin{axis}[
  width = 10.6cm, height = 6.6cm,
  axis lines = middle, xlabel = {$x$}, ylabel = {$y$},
  xmin = -3.4, xmax = 4.4, ymin = -3.4, ymax = 4.4,
  xtick = {-3,-2,-1,1,2,3,4}, ytick = {-3,-2,-1,1,2,3,4},
  axis line style = {bmtGris},
  tick label style = {font=\scriptsize, color=bmtGris},
  grid = both, grid style = {bmtGrisClair, line width=0.3pt},
  clip = true, restrict y to domain = -6:7,
]
  \addplot[bmtIndigo, line width=1.5pt, domain=-3.3:1.45, samples=180]
    {exp(x)};
  \addplot[bmtRaphia, line width=1.5pt, domain=0.05:4.3, samples=200]
    {ln(x)};
  \addplot[bmtGris, line width=0.8pt, dash pattern=on 2.5pt off 2pt,
           domain=-3.3:4.3] {x};
  \node[bmtIndigo, font=\scriptsize, anchor=south west]
    at (axis cs:1.1,3.3) {$y = e^{x}$};
  \node[bmtRaphia, font=\scriptsize, anchor=north west]
    at (axis cs:3.3,1.05) {$y = \ln x$};
  \node[bmtGris, font=\scriptsize, anchor=south east]
    at (axis cs:3.9,3.75) {$y = x$};
\end{axis}
\end{tikzpicture}
\end{center}

Les deux courbes sont symétriques par rapport à la droite $y = x$ : c'est la
traduction graphique de « l'une défait ce que l'autre fait ». On y lit aussi
que l'exponentielle reste toujours au-dessus de l'axe des abscisses, qui est
son asymptote horizontale en $-\infty$.
\end{visualisation}

% ===========================================================================
\section{Équations et inéquations}

\begin{theoreme}[les trois équations du programme]
Pour $a>0$ :
\[
  e^{u(x)} = a \daccord u(x) = \ln a, \qquad
  e^{u(x)} = e^{v(x)} \daccord u(x) = v(x).
\]
L'équation $a\,e^{2x}+b\,e^{x}+c = 0$ se résout en posant $X = e^{x}$, avec
la contrainte $X>0$ : on obtient un trinôme du second degré, et l'on ne
retient que les racines strictement positives.
\end{theoreme}

\begin{propriete}[inéquations]
L'exponentielle étant strictement croissante,
$e^{u(x)} < e^{v(x)} \daccord u(x) < v(x)$ : le sens de l'inégalité est
conservé.
\end{propriete}

\begin{avertissement}
$e^{u(x)} = a$ n'a \textbf{aucune} solution si $a \leq 0$ : une exponentielle
ne peut être ni nulle ni négative. C'est le contrôle à faire avant toute
autre chose, et c'est aussi ce qui fait rejeter les racines négatives dans un
changement d'inconnue.
\end{avertissement}

\begin{exemple}[le changement d'inconnue]
Résolvons $e^{2x}-4e^{x}+3 = 0$.

Posons $X = e^{x}$, avec $X>0$. L'équation devient $X^{2}-4X+3 = 0$, qui se
factorise en $(X-1)(X-3) = 0$ : $X = 1$ ou $X = 3$. Les deux sont
strictement positives, donc
\[
  e^{x} = 1 \;\text{ou}\; e^{x} = 3
  \quad\daccord\quad x = 0 \;\text{ou}\; x = \ln 3 .
\]
\end{exemple}

\begin{entrainement}[équations et inéquations]
\exo[\niveau{1}] $e^{x} = 5$

\exo[\niveau{1}] $e^{2x-1} = 1$

\exo[\niveau{2}] $e^{x^{2}} = e^{3x-2}$

\exo[\niveau{2}] $e^{x} = -2$

\exo[\niveau{2}] $e^{2x}-3e^{x}+2 = 0$

\exo[\niveau{3}] $e^{2x}-e^{x}-6 = 0$

\exo[\niveau{3}] $e^{x} < 4$

\exo[\niveau{3}] $\left(e^{x}-1\right)\left(e^{x}-3\right) \geq 0$
\end{entrainement}

\begin{corrige}[équations et inéquations]
\textbf{1.} $x = \ln5$.
\textbf{2.} $2x-1 = 0$ donc $x = \frac12$.
\textbf{3.} $x^{2} = 3x-2$, soit $x^{2}-3x+2 = 0$ : $x = 1$ ou $x = 2$.
\textbf{4.} Aucune solution : une exponentielle est strictement positive.
\textbf{5.} $X^{2}-3X+2 = 0$ donne $X = 1$ ou $X = 2$, d'où $x = 0$ ou
$x = \ln2$.
\textbf{6.} $X^{2}-X-6 = 0$ donne $X = -2$ (rejetée) ou $X = 3$ : $x = \ln3$.
\textbf{7.} $x<\ln4$, soit $\mathcal S = \intervalleo{-\infty}{\ln4}$.
\textbf{8.} Le produit est positif quand $e^{x} \leq 1$ ou $e^{x} \geq 3$,
soit $\mathcal S = \intervalleof{-\infty}{0}\cup\intervallefo{\ln3}{+\infty}$.
\end{corrige}

% ===========================================================================
\section{Dériver et étudier}

\begin{propriete}[dérivée d'une exponentielle composée]
Si $u$ est dérivable sur un intervalle $I$, alors $x \mapsto e^{u(x)}$ est
dérivable sur $I$ et
\[
  \left(e^{u}\right)'(x) = u'(x)\,e^{u(x)} .
\]
\end{propriete}

\begin{remarque}
Comme $e^{u(x)} > 0$, le signe de $\left(e^{u}\right)'$ est
\textbf{exactement} celui de $u'$. C'est ce qui rend l'étude de ces fonctions
si commode : le facteur exponentiel ne change jamais le signe, il ne fait que
l'amplifier.
\end{remarque}

\begin{exemple}[l'étude type du baccalauréat]
Soit $f(x) = x+2+e^{-x}$ sur $\R$.

\emph{Limites.} En $+\infty$, $e^{-x} \to 0$ et $x+2 \to +\infty$ : la
limite est $+\infty$. En $-\infty$, $e^{-x} \to +\infty$ : la limite est
$+\infty$ également.

\emph{Dérivée.} $f'(x) = 1-e^{-x} = \dfrac{e^{x}-1}{e^{x}}$. Le dénominateur
est positif : le signe est celui de $e^{x}-1$, négatif pour $x<0$, nul en
$0$, positif pour $x>0$.

\emph{Variations.} $f$ décroît sur $\R^{-}$, croît sur $\R^{+}$, minimum
$f(0) = 3$.

\emph{Asymptote oblique.} $f(x)-(x+2) = e^{-x} \to 0$ en $+\infty$ : la
droite $y = x+2$ est asymptote oblique en $+\infty$, et la courbe reste
au-dessus puisque $e^{-x}>0$.
\end{exemple}

\begin{entrainement}[dérivées et études]
\exo[\niveau{1}] Dériver $f(x) = e^{3x}$.

\exo[\niveau{1}] Dériver $g(x) = e^{-x}+x$.

\exo[\niveau{2}] Dériver $h(x) = x\,e^{x}$ et étudier le signe de $h'$.

\exo[\niveau{2}] Dériver $k(x) = \dfrac{e^{x}}{e^{x}+1}$ et donner son
signe.

\exo[\niveau{2}] $f(x) = e^{x}-x$ : dresser le tableau de variation sur
$\R$.

\exo[\niveau{3}] $f(x) = 4x^{2}e^{2x}$ : montrer que
$f'(x) = 8x(x+1)e^{2x}$, puis étudier le signe de $f'$.

\exo[\niveau{3}] $f(x) = e^{2x}-4e^{x}+3$ : montrer que
$f'(x) = 2e^{x}\left(e^{x}-2\right)$ et en déduire les variations.
\end{entrainement}

\begin{corrige}[dérivées et études]
\textbf{1.} $3e^{3x}$.
\textbf{2.} $1-e^{-x}$.
\textbf{3.} $h'(x) = (x+1)e^{x}$, du signe de $x+1$.
\textbf{4.} $k'(x) = \frac{e^{x}\left(e^{x}+1\right)-e^{x}\times e^{x}}
{\left(e^{x}+1\right)^{2}} = \frac{e^{x}}{\left(e^{x}+1\right)^{2}} > 0$.
\textbf{5.} $f'(x) = e^{x}-1$ : $f$ décroît sur $\R^{-}$, croît sur $\R^{+}$,
minimum $f(0) = 1$.
\textbf{6.} $f'(x) = 8xe^{2x}+4x^{2}\times2e^{2x} = 8x\left(1+x\right)e^{2x}$ ;
comme $e^{2x}>0$, le signe est celui de $x(x+1)$ : positif à l'extérieur de
$\intervalle{-1}{0}$, négatif à l'intérieur.
\textbf{7.} $f'(x) = 2e^{2x}-4e^{x} = 2e^{x}\left(e^{x}-2\right)$, du signe
de $e^{x}-2$ : négatif pour $x<\ln2$, positif après ; minimum
$f(\ln2) = 4-8+3 = -1$.
\end{corrige}

\begin{qcm}
\exo $e^{x} = 0$ admet :
\textbf{a)} une solution \quad \textbf{b)} deux solutions \quad
\textbf{c)} aucune solution \quad \textbf{d)} une infinité \reponseqcm{c}

\exo $\left(e^{-2x}\right)'$ vaut :
\textbf{a)} $e^{-2x}$ \quad \textbf{b)} $-2e^{-2x}$ \quad
\textbf{c)} $2e^{-2x}$ \quad \textbf{d)} $-2xe^{-2x-1}$ \reponseqcm{b}

\exo $\limminf e^{x}$ vaut :
\textbf{a)} $-\infty$ \quad \textbf{b)} $0$ \quad \textbf{c)} $1$ \quad
\textbf{d)} $+\infty$ \reponseqcm{b}

\exo $e^{a}\times e^{b}$ vaut :
\textbf{a)} $e^{ab}$ \quad \textbf{b)} $e^{a+b}$ \quad \textbf{c)}
$2e^{a+b}$ \quad \textbf{d)} $e^{a}+e^{b}$ \reponseqcm{b}

\exo L'équation $e^{2x}-e^{x}-2 = 0$ a pour solution :
\textbf{a)} $x = \ln2$ \quad \textbf{b)} $x = \ln2$ et $x = \ln(-1)$ \quad
\textbf{c)} $x = 2$ \quad \textbf{d)} aucune \reponseqcm{a}
\end{qcm}

% ===========================================================================
\begin{annales}
Comme pour le logarithme, le fonds est ici entièrement malgache : la moitié
des problèmes de la série~A repose sur une exponentielle. Les questions de
primitive et d'aire sont traitées au chapitre 9.

\exoann{Baccalauréat, Madagascar, série A, session 2021 — problème}
Soit $f(x) = x+2+e^{-x}$, de courbe $(\mathcal C)$ dans un repère orthonormé
d'unité $1$ cm.
\begin{enumerate}[label=\textbf{\alph*)}, leftmargin=*, itemsep=0.15em]
  \item Déterminer l'ensemble de définition de $f$.
  \item Montrer que, pour tout $x \neq 0$,
        $f(x) = x\left(1+\dfrac2x+\dfrac{1}{xe^{x}}\right)$ ; sachant que
        $\displaystyle\lim_{x \to -\infty} xe^{x} = 0^{-}$, calculer
        $\limminf f(x)$, puis $\limpinf f(x)$.
  \item Montrer que $f'(x) = \dfrac{e^{x}-1}{e^{x}}$.
  \item Étudier le signe de $e^{x}-1$ sur $\R$, puis dresser le tableau de
        variation de $f$.
  \item Montrer que la droite $(D) : y = x+2$ est asymptote oblique à
        $(\mathcal C)$.
  \item Calculer $f(-2)$ et $f(-1)$ à $0{,}01$ près, puis tracer la partie
        de $(\mathcal C)$ sur $\intervallefo{-2}{+\infty}$ en précisant son
        asymptote oblique.
\end{enumerate}

\exoann{Baccalauréat, Madagascar, série A, session 2014 — problème}
Soit $f(x) = 1-2x+e^{x}$, de courbe $(C)$ dans un repère orthonormé d'unité
$1$ cm.
\begin{enumerate}[label=\textbf{\alph*)}, leftmargin=*, itemsep=0.15em]
  \item Calculer $\limminf f(x)$.
  \item Vérifier que, pour tout $x \neq 0$,
        $f(x) = x\left(\dfrac1x-2+\dfrac{e^{x}}{x}\right)$ ; sachant que
        $\limpinf \dfrac{e^{x}}{x} = +\infty$, calculer $\limpinf f(x)$.
  \item Calculer $f'(x)$ et dresser le tableau de variation de $f$.
  \item Déterminer les coordonnées du point $A$, intersection de $(C)$ avec
        l'axe des ordonnées, puis écrire l'équation de la tangente $(T)$ à
        $(C)$ en $A$.
  \item Calculer $\displaystyle\lim_{x \to -\infty}
        \left[f(x)-(-2x+1)\right]$ ; que signifie ce résultat pour $(C)$ ?
  \item Calculer la valeur exacte de $f(1)$ et celle de $f(2)$, puis tracer
        $(T)$ et $(C)$ sur $\intervalleof{-\infty}{2}$.
\end{enumerate}

\exoann{Baccalauréat, Madagascar, série A, session 2013 — problème}
Soit $f(x) = \dfrac{4e^{x}}{e^{x}+1}$ sur $\intervallefo{0}{+\infty}$, de
courbe $(\mathscr C)$ d'unité graphique $1$ cm.
\begin{enumerate}[label=\textbf{\alph*)}, leftmargin=*, itemsep=0.15em]
  \item Vérifier que, pour tout $x \geq 0$,
        $f(x) = 4\left(1-\dfrac{1}{e^{x}+1}\right)$, puis en déduire
        $\limpinf f(x)$ et l'interpréter graphiquement.
  \item Prouver que, pour tout $x \geq 0$,
        $f'(x) = \dfrac{4e^{x}}{\left(e^{x}+1\right)^{2}}$, puis dresser le
        tableau de variation.
  \item Calculer $f(0)$, $f(1)$ et $f(2)$ à $10^{-2}$ près, écrire
        l'équation de la tangente $(T)$ au point d'abscisse $0$, et tracer
        $(T)$ et $(\mathscr C)$ avec leur asymptote.
\end{enumerate}

\exoann{Baccalauréat, Madagascar, série A, session 2017 — problème}
Soit $f(x) = \dfrac{e^{x}}{e^{x}+1}+2$ sur $\R$, de courbe $(\mathscr C)$
d'unité $2$ cm.
\begin{enumerate}[label=\textbf{\alph*)}, leftmargin=*, itemsep=0.15em]
  \item Calculer $\limminf f(x)$ et interpréter graphiquement le résultat.
  \item Démontrer que, pour tout réel $x$,
        $f(x) = 3-\dfrac{1}{e^{x}+1}$, puis en déduire $\limpinf f(x)$.
  \item Vérifier que $f'(x) = \dfrac{e^{x}}{\left(e^{x}+1\right)^{2}}$ et
        dresser le tableau de variation.
  \item Montrer que le point $I\left(0\,;\,\frac52\right)$ est centre de
        symétrie de $(\mathscr C)$, écrire l'équation de la tangente $(T)$ au
        point d'abscisse $0$, puis calculer $f(-1)$, $f(1)$ et $f(2)$ à
        $0{,}1$ près et tracer.
\end{enumerate}

\exoann{Baccalauréat, Madagascar, série A, session 2005 — problème
(extrait)}
Soit $f(x) = 4x^{2}e^{2x}$ sur $\R$, de courbe $(C)$ d'unité $4$ cm.
\begin{enumerate}[label=\textbf{\alph*)}, leftmargin=*, itemsep=0.15em]
  \item Calculer $\limpinf f(x)$ et $\limminf f(x)$ (on remarquera que
        $f(x) = \left[2xe^{x}\right]^{2}$ et l'on admettra que
        $\limminf xe^{x} = 0$), puis interpréter graphiquement.
  \item Montrer que $f'(x) = 8x(x+1)e^{2x}$ et en déduire son signe.
  \item Dresser le tableau de variation, puis écrire l'équation de la
        tangente à $(C)$ au point d'abscisse $-1$.
\end{enumerate}

\exoann{Baccalauréat, Madagascar, série A, session 2001 — problème
(extrait)}
Soit $f(x) = (1-x)e^{x}-1$ sur $\R$.
\begin{enumerate}[label=\textbf{\alph*)}, leftmargin=*, itemsep=0.15em]
  \item Calculer $\limpinf f(x)$ ; sachant que $\limminf xe^{x} = 0$,
        montrer que $\limminf f(x) = -1$ et interpréter graphiquement.
  \item Calculer $f'(x)$, étudier son signe et dresser le tableau de
        variation.
  \item Déterminer les coordonnées du point $A$, intersection de la courbe
        avec la droite $(D) : y = -1$, puis écrire l'équation de la tangente
        $(T)$ au point d'abscisse $1$.
\end{enumerate}

\exoann{Baccalauréat blanc, Lycée Philibert Tsiranana, série A, 2023 —
problème (extrait)}
Soit $f(x) = e^{2x}-4e^{x}+3$ sur $\R$.
\begin{enumerate}[label=\textbf{\alph*)}, leftmargin=*, itemsep=0.15em]
  \item Calculer $\limminf f(x)$ et interpréter le résultat.
  \item En écrivant $f(x) = e^{x}\left(e^{x}-4\right)+3$, calculer
        $\limpinf f(x)$ puis $\limpinf \dfrac{f(x)}{x}$ ; en déduire la
        nature de la branche infinie en $+\infty$.
  \item Montrer que $f'(x) = 2e^{x}\left(e^{x}-2\right)$, étudier son signe
        et dresser le tableau de variation.
  \item Vérifier que $f(x) = \left(e^{x}-1\right)\left(e^{x}-3\right)$,
        résoudre $f(x) = 0$, et en déduire les points d'intersection de la
        courbe avec l'axe des abscisses.
\end{enumerate}
\end{annales}

\begin{corrige}[annales]
\textbf{1.} $\Df = \R$. En $-\infty$ : $e^{-x} \to +\infty$, donc
$f \to +\infty$ ; en $+\infty$ : $e^{-x} \to 0$ et $x+2 \to +\infty$, donc
$f \to +\infty$. La dérivée vaut $1-e^{-x} = \frac{e^{x}-1}{e^{x}}$, du
signe de $e^{x}-1$, c'est-à-dire de $x$ : minimum $f(0) = 3$. Enfin
$f(x)-(x+2) = e^{-x} \to 0$ en $+\infty$ : $(D)$ est asymptote oblique, et
la courbe reste au-dessus. $f(-2) = 7{,}39$ et $f(-1) = 3{,}72$.

\textbf{2.} En $-\infty$, $e^{x} \to 0$ et $-2x \to +\infty$ :
$f \to +\infty$. La factorisation donne une parenthèse de limite $+\infty$,
donc $f \to +\infty$ en $+\infty$. $f'(x) = -2+e^{x}$, qui s'annule pour
$x = \ln2$ : minimum $f(\ln2) = 1-2\ln2+2 = 3-2\ln2 \approx 1{,}61$. Le
point $A$ est $(0\,;\,2)$ et $f'(0) = -1$, d'où $(T) : y = -x+2$. Enfin
$f(x)-(-2x+1) = e^{x} \to 0$ en $-\infty$ : la droite $y = -2x+1$ est
asymptote oblique en $-\infty$. Valeurs exactes : $f(1) = e-1$ et
$f(2) = e^{2}-3$.

\textbf{3.} $4\left(1-\frac{1}{e^{x}+1}\right)
= \frac{4\left(e^{x}+1\right)-4}{e^{x}+1} = \frac{4e^{x}}{e^{x}+1}$
\checkmark. Comme $\frac{1}{e^{x}+1} \to 0$, $f \to 4$ : la droite $y = 4$
est asymptote horizontale. La dérivée est strictement positive : $f$ est
croissante, de $f(0) = 2$ vers $4$. Valeurs : $f(1) \approx 2{,}92$,
$f(2) \approx 3{,}52$ ; $f'(0) = 1$, donc $(T) : y = x+2$.

\textbf{4.} En $-\infty$, $e^{x} \to 0$ donc $f \to 0+2 = 2$ : asymptote
horizontale $y = 2$. L'écriture $3-\frac{1}{e^{x}+1}$ se vérifie par
réduction au même dénominateur, et donne $f \to 3$ en $+\infty$ : seconde
asymptote horizontale. $f'>0$ : la fonction croît de $2$ à $3$. Pour la
symétrie, $f(-x)+f(x) = 5 = 2\times\frac52$. Enfin $f(0) = \frac52$ et
$f'(0) = \frac14$, d'où $(T) : y = \frac14x+\frac52$ ;
$f(-1) \approx 2{,}3$, $f(1) \approx 2{,}7$, $f(2) \approx 2{,}9$.

\textbf{5.} $f(x) = \left[2xe^{x}\right]^{2} \geq 0$ ; en $-\infty$, le
crochet tend vers $0$ donc $f \to 0$ : l'axe des abscisses est asymptote
horizontale. En $+\infty$, $f \to +\infty$. La dérivée
$8x(x+1)e^{2x}$ est du signe de $x(x+1)$ : maximum local
$f(-1) = 4e^{-2} \approx 0{,}54$ en $-1$, minimum $f(0) = 0$. Comme
$f'(-1) = 0$, la tangente en $-1$ est horizontale : $y = 4e^{-2}$.

\textbf{6.} En $+\infty$, $(1-x) \to -\infty$ et $e^{x} \to +\infty$ :
$f \to -\infty$. En $-\infty$, $f(x) = e^{x}-xe^{x}-1 \to 0-0-1 = -1$ : la
droite $y = -1$ est asymptote horizontale. $f'(x) = -e^{x}+(1-x)e^{x}
= -xe^{x}$, du signe de $-x$ : maximum $f(0) = 0$. Le point $A$ vérifie
$(1-x)e^{x}-1 = -1$, soit $(1-x)e^{x} = 0$, donc $x = 1$ : $A(1\,;\,-1)$.
Comme $f'(1) = -e$, la tangente en $1$ est $y = -e(x-1)-1$.

\textbf{7.} En $-\infty$, $e^{x} \to 0$ donc $f \to 3$ : asymptote
horizontale $y = 3$. En $+\infty$, $e^{x}\left(e^{x}-4\right) \to +\infty$
donc $f \to +\infty$, et $\frac{f(x)}{x} \to +\infty$ : branche parabolique
de direction $(Oy)$. $f'(x) = 2e^{x}\left(e^{x}-2\right)$, du signe de
$e^{x}-2$ : minimum $f(\ln2) = 4-8+3 = -1$. Enfin
$\left(e^{x}-1\right)\left(e^{x}-3\right) = e^{2x}-4e^{x}+3$ \checkmark, et
$f(x) = 0$ donne $x = 0$ ou $x = \ln3$ : les points sont $(0\,;\,0)$ et
$(\ln3\,;\,0)$.
\end{corrige}

% ===========================================================================
\begin{vraifaux}
\exo $e^{x}$ peut être négatif si $x$ est négatif.

\exo $e^{x+y} = e^{x}+e^{y}$.

\exo La courbe de $\exp$ admet l'axe des abscisses pour asymptote en
$-\infty$.

\exo La dérivée de $e^{-x}$ est $e^{-x}$.

\exo Si $e^{a} = e^{b}$ alors $a = b$.
\end{vraifaux}

\begin{corrige}[vrai ou faux]
\textbf{1.} Faux — $e^{x}>0$ pour tout réel $x$ ; pour $x$ négatif, $e^{x}$
est simplement \emph{petit}.
\textbf{2.} Faux — c'est $e^{x}\times e^{y}$.
\textbf{3.} Vrai — $\limminf e^{x} = 0$.
\textbf{4.} Faux — elle vaut $-e^{-x}$.
\textbf{5.} Vrai — l'exponentielle est strictement croissante.
\end{corrige}

% ===========================================================================
\begin{situation}[la population d'une commune]
La population d'une commune, en milliers d'habitants,
est modélisée par
\[
  P(t) = 12\,e^{0{,}021\,t},
\]
où $t$ désigne le nombre d'années écoulées depuis le dernier recensement.

\begin{enumerate}[label=\textbf{\arabic*.}, leftmargin=*, itemsep=0.3em]
  \item Quelle était la population au moment du recensement ?
  \item Calculer la population prévue dix ans plus tard, au millier près.
  \item Montrer que la population double au bout de
        $t = \dfrac{\ln 2}{0{,}021}$ années, puis en donner une valeur
        approchée. On donne $\ln 2 \approx 0{,}693$.
  \item Étudier le sens de variation de $P$ et dire ce que le modèle suppose
        implicitement.
\end{enumerate}
\end{situation}

\begin{corrige}[la population d'une commune]
\textbf{1.} $P(0) = 12$, soit $12\,000$ habitants.

\textbf{2.} $P(10) = 12e^{0{,}21} \approx 12\times1{,}234 \approx 14{,}8$,
soit environ $15\,000$ habitants.

\textbf{3.} On cherche $t$ tel que $12e^{0{,}021t} = 24$, soit
$e^{0{,}021t} = 2$, d'où $0{,}021\,t = \ln2$ et
$t = \frac{\ln2}{0{,}021} \approx 33$ ans.

\textbf{4.} $P'(t) = 12\times0{,}021\,e^{0{,}021t} > 0$ : la population
croît sans cesse, et de plus en plus vite. Le modèle suppose un taux de
croissance constant et sans limite — ce qui ne peut pas rester vrai
indéfiniment, faute de terres et de ressources.
\end{corrige}

% ===========================================================================
\begin{jemeteste}
Trente minutes.

\begin{enumerate}[label=\textbf{\arabic*.}, leftmargin=*, itemsep=0.3em]
  \item Simplifier $e^{3}\times e^{-5}$ et $\dfrac{e^{2x}}{e^{x}}$.
  \item Résoudre $e^{3x-2} = e^{x+4}$.
  \item Résoudre $e^{2x}-5e^{x}+6 = 0$.
  \item Dériver $f(x) = (2x-1)e^{x}$.
  \item $g(x) = e^{x}-x-1$ : dresser le tableau de variation et en déduire
        que $e^{x} \geq x+1$ pour tout réel $x$.
\end{enumerate}
\end{jemeteste}

\begin{corrige}[je me teste]
\textbf{1.} $e^{-2}$ ; et $e^{x}$.
\textbf{2.} $3x-2 = x+4$ donne $x = 3$.
\textbf{3.} $X = 2$ ou $X = 3$, d'où $x = \ln2$ ou $x = \ln3$.
\textbf{4.} $f'(x) = 2e^{x}+(2x-1)e^{x} = (2x+1)e^{x}$.
\textbf{5.} $g'(x) = e^{x}-1$, du signe de $x$ : minimum $g(0) = 0$. Le
minimum étant nul, $g(x) \geq 0$ pour tout $x$, c'est-à-dire
$e^{x} \geq x+1$.
\end{corrige}

% ===========================================================================
\begin{commeaubac}
\bmtsource{Baccalauréat, Madagascar, série A, session 2003 — problème
\textbullet\ 10 points}
Soit $f$ la fonction définie par $f(x) = 1-2x+e^{x}$. On note $(C)$ sa
courbe représentative dans un repère orthonormé
$\left(O\,;\,\veci\,,\,\vecj\right)$ d'unité $1$ cm.

\begin{enumerate}[label=\textbf{\arabic*.}, leftmargin=*, itemsep=0.25em]
  \item \begin{enumerate}[label=\textbf{\alph*)}, leftmargin=1.4em, itemsep=0.15em]
      \item Déterminer l'ensemble de définition de $f$. \bareme{0,50}
      \item Calculer $\limminf f(x)$. \bareme{0,50}
      \item En remarquant que, pour tout $x>0$,
            $f(x) = x\left(\dfrac1x-2+\dfrac{e^{x}}{x}\right)$, calculer
            $\limpinf f(x)$. On donne $\limpinf \dfrac{e^{x}}{x} = +\infty$.
            \bareme{0,50}
    \end{enumerate}
  \item \begin{enumerate}[label=\textbf{\alph*)}, leftmargin=1.4em, itemsep=0.15em]
      \item Calculer $f'(x)$. \bareme{0,75}
      \item En déduire le tableau de variation de $f$. \bareme{1}
    \end{enumerate}
  \item \begin{enumerate}[label=\textbf{\alph*)}, leftmargin=1.4em, itemsep=0.15em]
      \item Déterminer les coordonnées du point $A$, intersection de la
            courbe $(C)$ avec l'axe des ordonnées. \bareme{0,75}
      \item Écrire l'équation de la tangente $(T)$ à $(C)$ au point $A$.
            \bareme{1}
    \end{enumerate}
  \item \begin{enumerate}[label=\textbf{\alph*)}, leftmargin=1.4em, itemsep=0.15em]
      \item Calculer $\displaystyle\lim_{x \to -\infty}
            \left[f(x)-(-2x+1)\right]$. Que peut-on en conclure ?
            \bareme{1}
      \item Étudier la branche infinie de $(C)$ lorsque $x$ tend vers
            $+\infty$. \bareme{0,75}
    \end{enumerate}
  \item Tracer $(C)$. \bareme{1,25}
\end{enumerate}

\begin{remarque}
La question 6 de ce sujet — primitive de $f$ puis aire du domaine limité par
$(C)$, l'axe des abscisses et les droites $x = 0$ et $x = \ln2$ — figure au
chapitre 9.
\end{remarque}
\end{commeaubac}

\begin{corrige}[comme au baccalauréat — Madagascar A 2003]
\textbf{1. a)} $f$ est une somme de fonctions définies sur $\R$ :
$\Df = \R$.
\textbf{b)} En $-\infty$, $e^{x} \to 0$ et $-2x \to +\infty$ :
$f \to +\infty$.
\textbf{c)} La parenthèse tend vers $+\infty$ et $x \to +\infty$ : donc
$f \to +\infty$.

\textbf{2. a)} $f'(x) = -2+e^{x}$.
\textbf{b)} $f'(x) = 0$ pour $e^{x} = 2$, soit $x = \ln2$ ; $f'$ est
négative avant, positive après. La fonction décroît sur
$\intervalleof{-\infty}{\ln2}$, croît sur $\intervallefo{\ln2}{+\infty}$, et
son minimum vaut $f(\ln2) = 1-2\ln2+2 = 3-2\ln2 \approx 1{,}61$.

\textbf{3. a)} L'axe des ordonnées a pour équation $x = 0$, et
$f(0) = 1-0+1 = 2$ : $A(0\,;\,2)$.
\textbf{b)} $f'(0) = -2+1 = -1$, d'où $(T) : y = -x+2$.

\textbf{4. a)} $f(x)-(-2x+1) = e^{x} \to 0$ en $-\infty$ : la droite
$y = -2x+1$ est asymptote oblique à $(C)$ en $-\infty$.
\textbf{b)} En $+\infty$, $\frac{f(x)}{x} = \frac1x-2+\frac{e^{x}}{x} \to
+\infty$ : la courbe admet une branche parabolique de direction $(Oy)$.

\textbf{5.} On trace l'asymptote oblique $y = -2x+1$, la tangente
$y = -x+2$, le sommet $\left(\ln2\,;\,3-2\ln2\right) \approx
(0{,}69\,;\,1{,}61)$, puis quelques points : $f(1) = e-1 \approx 1{,}72$ et
$f(2) = e^{2}-3 \approx 4{,}39$. La courbe descend depuis le haut à gauche,
touche son minimum, puis remonte de plus en plus vite.
\end{corrige}

\begin{notepourlaclasse}
Quatre heures suffisent si le chapitre 7 a été solidement fait : tout se
déduit du logarithme, et il n'y a presque rien à mémoriser en plus. Le seul
point réellement nouveau est le changement d'inconnue $X = e^{x}$, avec sa
contrainte $X>0$ — à faire écrire, systématiquement, avant la résolution.

L'inégalité $e^{x} \geq x+1$ du bloc « je me teste » vaut d'être commentée :
c'est la même que $\ln x \leq x-1$, lue de l'autre côté de la droite
$y = x$. Les élèves y voient d'un coup à quoi sert la symétrie des deux
courbes.

Deux sujets suffisent en classe — 2003 et 2021 — le reste étant à donner en
travail personnel. Notez que 2003 et 2014 posent, à onze ans d'intervalle,
exactement la même fonction : la remarque frappe toujours les élèves et
justifie à elle seule le travail sur annales.
\end{notepourlaclasse}
